\documentclass[12pt]{article}
\newcommand{\DocShort}{Curva logística}
\input{preamble_population.tex}
\begin{document}
\doccover{Método de la curva logística}{2026-05}
\handoutintro{Proyectar la población mediante una curva en S con saturación finita.}{Ciudades grandes donde el crecimiento ya muestra señales de acercamiento a un límite práctico.}

\section{Datos ejemplo}
Se necesitan tres poblaciones tomadas en orden cronológico y separadas por intervalos iguales:

\begin{center}
\begin{tabular}{lccc}
\toprule
Año & 2001 & 2011 & 2021 \\
\midrule
Población & 36,300 & 76,400 & 118,700 \\
\bottomrule
\end{tabular}
\end{center}

En este documento:
\[
P_0=\num{36300}\ \text{en 2001},\qquad P_1=\num{76400}\ \text{en 2011},\qquad P_2=\num{118700}\ \text{en 2021}
\]
Los intervalos son iguales y valen $10$ años, de modo que $t_0=0$, $t_1=10$ años y $t_2=20$ años.

\section{Ecuaciones mínimas}
La ecuación logística se escribe como:
\begin{equation}
P(t)=\frac{P_s}{1+me^{nt}}
\end{equation}
donde $P(t)$ es la población al tiempo $t$, $P_s$ es la población de saturación, $m$ y $n$ son constantes del modelo y $t$ es el tiempo medido en años desde el primer dato.

La población de saturación se estima con:
\begin{equation}
P_s=\frac{2P_0P_1P_2-P_1^2(P_0+P_2)}{P_0P_2-P_1^2}
\end{equation}

Luego se calculan:
\begin{equation}
m=\frac{P_s-P_0}{P_0}
\end{equation}

\begin{equation}
n=\frac{2.3}{t_1}\log_{10}\left(\frac{P_0(P_s-P_1)}{P_1(P_s-P_0)}\right)
\end{equation}

\section{Procedimiento}
\begin{enumerate}
\item Ordenar cronológicamente las tres poblaciones como $P_0$, $P_1$ y $P_2$.
\item Verificar que los intervalos entre observaciones sean iguales.
\item Calcular la población de saturación $P_s$.
\item Calcular las constantes $m$ y $n$.
\item Sustituir en la ecuación logística y evaluar el año deseado.
\end{enumerate}

\section{Ejemplo resuelto}
\begin{center}
\begin{tikzpicture}
\begin{axis}[
width=0.88\textwidth,
height=6.6cm,
xlabel={Año},
ylabel={Población},
xmin=2000, xmax=2040,
ymin=0, ymax=170000,
grid=both,
major grid style={gray!25},
minor grid style={gray!10},
xtick={2000,2005,2010,2015,2020,2025,2030,2035,2040},
ytick={0,20000,40000,60000,80000,100000,120000,140000,160000},
scaled y ticks=false,
xticklabel style={/pgf/number format/1000 sep={}}
]
\addplot[smooth, thick, color=teal, domain=2001:2040, samples=120]
{161203/(1 + 3.44*exp(-0.113*(x-2001)))};
\addplot[only marks, mark=*, mark size=2.5pt, color=medblue]
coordinates {(2001,36300) (2011,76400) (2021,118700)};
\end{axis}
\end{tikzpicture}
\end{center}

La gráfica muestra que los datos son consistentes con una trayectoria de tipo sigmoidal o curva S.

\[
P_s=\frac{2(\num{36300})(\num{76400})(\num{118700})-\num{76400}^2(\num{36300}+\num{118700})}{(\num{36300})(\num{118700})-\num{76400}^2}=\num{161203}
\]

\[
m=\frac{\num{161203}-\num{36300}}{\num{36300}}=3.44
\]

\[
n=\frac{2.3}{10}\log_{10}\left(\frac{\num{36300}(\num{161203}-\num{76400})}{\num{76400}(\num{161203}-\num{36300})}\right)=-0.113
\]

La ecuación queda:
\[
P(t)=\frac{\num{161203}}{1+3.44e^{-0.113t}}
\]

Para 2031, desde 2001 han transcurrido $t=30$ años:
\[
P_{2031}=\frac{\num{161203}}{1+3.44e^{-0.113(30)}}=\num{144454}
\]

\resultbox{La curva logística proyecta 144,454 habitantes para el año 2031.}

\section{Teoría breve}
La curva logística modela crecimiento rápido al inicio, casi lineal en la zona media y desacelerado cerca de la saturación. Por eso es un modelo más estructurado que los métodos puramente extrapolativos.

\section{Observaciones de uso}
\begin{itemize}
\item Es sensible a la selección de $P_0$, $P_1$ y $P_2$.
\item Los intervalos entre observaciones deben ser iguales.
\item El valor de saturación debe interpretarse como límite práctico, no como una constante física exacta.
\end{itemize}
\end{document}
